\documentclass[a4paper,14pt]{extarticle}
\usepackage[T2A]{fontenc}
\usepackage[utf8]{inputenc}
\usepackage[russian]{babel}
\usepackage{tempora,setspace,indentfirst}
\usepackage[left=30mm,right=15mm,top=20mm,bottom=20mm]{geometry}
\onehalfspacing
\setlength{\parindent}{1.25cm}
\begin{document}
\noindent\textit{Учебный пример. Результаты и условия проверки вымышлены.}
\tableofcontents
\clearpage
\section*{\centering\normalsize ЗАКЛЮЧЕНИЕ}
\addcontentsline{toc}{section}{Заключение}
В работе разработан прототип учета заявок, обеспечивающий регистрацию,
изменение статуса и поиск записей. Полученное решение позволяет представить
основные операции учета в едином программном интерфейсе.

При анализе задачи выделены три группы функциональных требований.
Спроектирована модель заявки с тремя статусами и допустимыми переходами между ними.
На ее основе реализованы функции регистрации, изменения статуса и поиска.

Функциональная проверка на учебном стенде включала десять сценариев:
четыре для регистрации, три для изменения статуса и три для поиска.
Во всех сценариях получен ожидаемый результат. Тем самым подтверждена работа
реализованных функций в проверенных условиях; цель разработки и функциональной
проверки прототипа достигнута в заданных границах.

Прототип может служить основой для последующей опытной эксплуатации.
Проверка проводилась с одним пользователем и учебными данными, поэтому ее результаты
не подтверждают производительность при одновременной работе сотрудников.
Дальнейшее развитие предполагает разграничение прав доступа и нагрузочные испытания.
Внедрение и сокращение времени обработки заявок в данной работе не оценивались.
\end{document}
